\begin{center}\begin{tikzpicture}[font=\small,>=stealth,line join=round]
\draw[gray!40,rounded corners=3pt] (0,0) rectangle (14.2,7.4);
\node[text width=12.6cm,align=center] at (7.1,6.75) {中心的高度不能丢：偶函数配 $(1,2)$ 的例子};

\begin{scope}[shift={(2.4,2.5)},x=1.55cm,y=.85cm]
\draw[->,gray] (-.8,0)--(5.3,0) node[right]{$x$};
\draw[->,gray] (0,-.3)--(0,3.7) node[above]{$y$};
\draw[gray,dashed] (-.5,2)--(5,2) node[right]{$y=2$};
\draw[blue,thick,domain=-.5:5,samples=200] plot(\x,{2+cos(deg(pi*\x/2))});
\fill[violet] (1,2) circle(2.7pt);\node[above right=4pt] at (1,2){$(1,2)$};
\foreach \u/\v in {0/3,2/1,4/3}{\fill[orange!85!black] (\u,\v) circle(2.3pt);\node[below=4pt] at (\u,0){$\u$};}
\end{scope}
\node at (7.1,1.35) {$f(0)=3,\quad f(2)=1\ne-3,\quad f(x+2)=4-f(x)$};

\node[text width=12.6cm,align=center] at (7.1,.65) {$f(x)=2+\cos(\pi x/2)$ 为偶函数；点对称使对应高度之和为 $4$，不是 $0$。};
\end{tikzpicture}\end{center}
